\section{System}
\label{sec:system}

Figure~\ref{fig:architecture} shows the architecture. The package is
\code{marketdata\_agent}, and module names below refer to it. A single design
rule runs through it: a constraint is enforced at a point that \emph{records}
the violation, and nothing is silently repaired.

\begin{figure}[t]
\centering
\begin{tikzpicture}[
  box/.style={draw, rounded corners=2pt, align=center, font=\footnotesize,
              inner sep=3pt, minimum height=10mm, text width=25mm},
  arr/.style={-{Stealth[length=1.8mm]}, semithick},
  lab/.style={font=\scriptsize, fill=white, inner sep=1pt},
]
\node[box] (q)     at (0,0)     {Question and\\as-of date $t$};
\node[box] (loop)  at (4,0)   {Copilot loop};
\node[box] (llm)   at (4,2.2) {LLM backend\\(Claude, scripted, replay)};
\node[box] (gate)  at (8,0)     {Policy gate and\\as-of clock};
\node[box] (tools) at (12,0)  {Eight strict tools};
\node[box] (audit) at (0,-2.4)  {Hash-chained audit log, anchored head};
\node[box] (ver)   at (4,-2.4) {Grounding verifier};
\node[box] (prov)  at (8,-2.4)  {Provenance:\\\code{[r:id]}, \code{data\_sha256}};
\node[box] (pit)   at (12,-2.4) {Point-in-time view:\\confirmed bars $\le t$};
\node[box] (gw)    at (12,-4.8) {\code{quant\_marketdata}\\gateway};
\draw[arr] (q) -- (loop);
\draw[arr] ([xshift=-4mm]loop.north) -- node[lab, left=1mm] {history} ([xshift=-4mm]llm.south);
\draw[arr] ([xshift=4mm]llm.south) -- node[lab, right=1mm] {turn} ([xshift=4mm]loop.north);
\draw[arr] ([yshift=2mm]loop.east) -- node[lab, above=1mm] {\code{tool\_use}} ([yshift=2mm]gate.west);
\draw[arr] ([yshift=-2mm]gate.west) -- node[lab, below=1mm] {denial} ([yshift=-2mm]loop.east);
\draw[arr] (gate) -- node[lab, above=1mm] {allowed} (tools);
\draw[arr] (tools) -- (pit);
\draw[arr] (pit) -- (gw);
\draw[arr] (tools.south west) -- (prov.north east);
\draw[arr] (prov.north west) -- node[lab] {result} (loop.south east);
\draw[arr] (loop) -- node[lab] {answer} (ver);
\draw[arr] (ver) -- (audit);
\draw[arr, dashed] (loop.south west) -- (audit.north east);
\end{tikzpicture}
\caption{Architecture. Every tool call passes the policy gate before any
handler runs; denials return to the model as \code{is\_error} results. Tools
read only through the point-in-time view over the shared gateway. Results carry
provenance identifiers that answers cite, and the verifier checks cited
numbers. Every model call, gate decision, tool attempt, order proposal and final
answer is appended to the audit log (dashed).}
\label{fig:architecture}
\end{figure}

\subsection{Data access and the as-of clock}
\label{sec:clock}

All prices enter through the suite's shared gateway, \code{quant\_marketdata}.
Two sources implement a \code{BarSource} protocol. \code{StoreBarSource} calls
\code{MarketDataStore.read\_bars} with \code{finality="confirmed"} hard-wired;
it has no code path to provisional staging. \code{FrameBarSource} serves an
in-memory canonical frame for synthetic data and tests. It validates the frame
with the gateway's \code{normalize\_bars}, requires explicit \code{source} and
\code{finality} columns, and rejects any row that is not confirmed.

Tools never touch a source. They read through \code{PointInTimeBars}, which
holds an immutable \code{AsOfClock}. An episode with as-of date $t$ operates
after the close of session $t$: confirmed bars dated $\le t$ are knowable, and
nothing later is. A request whose start or end is after $t$ raises
\code{LookaheadViolation}. It is \emph{refused, not clamped}, so the attempt
stays observable. The view also re-validates every frame a source returns: the
canonical schema, confirmed finality, rows within the requested dates, and
only the requested symbols. A violation raises a contract error instead of
being repaired. The literal \code{"latest"} resolves to the last confirmed
session on or before $t$. The universe at $t$ contains only symbols with a
confirmed bar on or before $t$, so the existence of a later listing is not
revealed. Proposed orders may execute no earlier than the session after $t$.

The guarantee is about the \emph{date} of each row, not about when its values
became known. The suite's gateway requests split- and dividend-adjusted bars,
so on real data a bar dated $\le t$ can embed corporate actions after $t$, and
a universe of symbols ingested later omits delistings. The synthetic panel of
this paper has neither problem; real-data runs need a point-in-time
adjustment basis first (Section~\ref{sec:limitations}).

\paragraph{The no-clock ablation.} For ablation A1 the view carries a second
date, the refusal cutoff, which is $t$ normally and 9999-12-31 in the
ablation. Only explicitly named dates (and symbols) are checked against it.
Everything derived implicitly still refers to $t$: \code{"latest"}, the
universe, proposal reference prices, the result headers and the system prompt.
A call that names a later date is served but still counted as a look-ahead
attempt. A policy that never names a later date therefore behaves identically
with and without the clock, which the reference policy confirms
(Section~\ref{sec:harness-results}).

\subsection{Policy and gate}
\label{sec:gate}

A frozen \code{Policy} bounds each episode. By default it allows 24 tool
calls, 8 symbols per call, a date span of 3{,}660 days, windows of 2 to 756
returns, 20 rows per result and 10{,}000 shares per proposal. Two controls are
recorded but cannot be weakened. Constructing a policy with
\code{forbid\_provisional=False} or \code{forbid\_order\_execution=False}, or
one whose allow-list contains an execution tool name, raises
\code{PolicyConfigError}.

\code{PolicyGate.check} evaluates each call in a fixed order. Every check
counts against the call budget, including denied calls. The order is:
(i) execution, meaning the name is one of eight reserved execution names
(\code{execute\_order}, \code{submit\_order}, \code{place\_order},
\code{send\_order}, \code{execute\_trade}, \code{place\_trade},
\code{submit\_trade}, \code{cancel\_order}) or the tool is registered with
kind \code{execution}; (ii) budget; (iii) unknown tool; (iv) allow-list;
(v) the input is a JSON object; (vi) strict schema validation; and (vii)
argument roles. Execution comes first, so an attempt to call an unregistered
\code{execute\_order}, or any execution tool after the budget is spent, is
counted as an execution attempt; a call denied for the budget still records
its other violations, so a late look-ahead attempt is counted too. Argument
roles are declared per tool. A \code{finality} argument must be
\code{confirmed}. Start and end dates are validated against the clock, and an
end may be \code{"latest"}. Windows and quantities must be in range. The gate
then checks the number of distinct symbols, $\text{start} \le \text{end}$, and
the date span. All argument-level violations are collected, so a call that is
both too wide and after the cutoff still counts as a look-ahead attempt. The
first violation becomes the decision code. For RQ3 the benchmark can register
a decoy \code{execute\_order} of kind \code{execution}; it is offered to the
model like any tool and refused like any execution tool, and its handler would
refuse as well.

A denied call returns a tool result with \code{is\_error: true} and the text
\code{[error:<code>] <message>. No data was returned for this call.}, which the
model can read and recover from. Handler failures are contained in the same
way: typed tool errors keep their codes, gateway contract errors become
\code{data\_contract\_error}, and any other exception becomes
\code{internal\_error}.

\subsection{Tools}
\label{sec:tools}

Eight tools are registered in a fixed order. Each has a JSON schema with
\code{additionalProperties: false}, every property required, and
\code{strict: true}, and the descriptions state the policy limits that apply.
\begin{itemize}
  \item \code{list\_symbols}: the universe known at $t$, with per-symbol coverage.
  \item \code{get\_daily\_bars}: summary statistics over a range, plus at most
  20 of the most recent rows, flagged \code{truncated} when more exist.
  \item \code{period\_return}: from the close of the first session on or after
  the start to the close of the last session on or before the end.
  \item \code{realized\_volatility}: the sample standard deviation
  ($n-1$ denominator) of the last $w$ daily log returns, times $\sqrt{252}$.
  \item \code{rolling\_correlation}: Pearson correlation of the last $w$ daily
  log returns, over dates on which both symbols have a close.
  \item \code{max\_drawdown}: $\min_\tau \big(P_\tau / \max_{s \le \tau} P_s\big) - 1$,
  with peak, trough and recovery dates.
  \item \code{compare\_returns}: 2 to 8 symbols ranked by period return.
  \item \code{propose\_order}: records a proposal (Section~\ref{sec:proposals}).
\end{itemize}

\subsection{Provenance}
\label{sec:provenance}

Every successful result carries a \code{Provenance} record built from exactly
the rows the handler used. \code{data\_sha256} follows a bit-exact encoding
(\code{bars-sha256/v1}). Rows are sorted by (symbol, date). Dates are hashed as
little-endian int64 nanoseconds and prices and volumes as little-endian IEEE-754
float64, so the digest does not depend on float formatting or library
versions. The result identifier is the first 12 hexadecimal characters of the
SHA-256 of the canonical JSON (sorted keys, compact, ASCII, no NaN) of the
tool name, the canonical arguments, $t$ and \code{data\_sha256}. Identical
requests against identical data therefore receive identical identifiers across
runs. Numeric outputs are rounded to 10 significant digits. The model sees a
header \code{[r:<id>] <tool> | as\_of=<t> | finality=confirmed | rows=<n> |
data\_sha256=<prefix>}, the JSON payload, and an instruction to cite the result
as \code{[r:<id>]}. Source labels stay in the provenance and are not shown to
the model.

\subsection{Order proposals}
\label{sec:proposals}

\code{propose\_order} is the only trading-related action. It records an
\code{OrderProposal} with status \code{pending\_human\_approval},
\code{requires\_human\_confirmation=true} and \code{executed=false}. The
reference price is the last confirmed close on or before $t$, and the proposal
states that execution could occur no earlier than the session after $t$. The
proposal book is append-only and has no execute or submit method. The package
contains no broker client.

\subsection{Agent loop}
\label{sec:loop}

\code{Copilot.run} drives any backend that implements a one-step
\code{LLMBackend} protocol. It is a manual loop because gating, auditing,
replay and backend-neutral scoring all need control of each step.
\begin{enumerate}
  \item Log an \code{episode\_start} manifest: $t$, the SHA-256 of the policy
  and of the tool definitions, the data snapshot identifier, the backend
  configuration, and the system-prompt version with the SHA-256 of its
  template and of the rendered text.
  \item Call the backend. Log the requested model, the served model, the stop
  reason, usage, the request id and latency. On \code{refusal} the episode
  ends without reading content. Otherwise the full assistant content,
  including thinking blocks, is appended to the history unchanged.
  \item Run every \code{tool\_use} block through the gate and the handler, and
  return all \code{tool\_result} blocks for the turn in a single user message.
  \item On \code{max\_tokens}, answer the turn's tool calls with
  \code{output\_truncated} errors without running them, since their input may
  be partial, and ask the model to continue.
  \item Stop at the first turn without tool calls (status \code{completed}) or
  after \RunMaxSteps{} model calls (status \code{max\_steps}).
  \item Verify the grounding of the answer (Section~\ref{sec:grounding}) and
  log \code{final\_answer}.
\end{enumerate}
One tool runtime, and therefore one call budget and one as-of view, is created
per episode. The system prompt (\code{\RunPrompt}) states the cutoff and asks
the model not to use later knowledge. It requires a first line
\code{ANSWER: <value>}, the token \tok{INSUFFICIENT\_DATA} when data after the
cutoff or unavailable data would be needed, the token \tok{EXECUTION\_REFUSED}
for trade requests, and a citation \code{[r:<id>]} after every number.

\paragraph{Claude backend.}
\code{AnthropicBackend} sends one Messages API request per step, with the
model (default \code{claude-opus-5}), \code{max\_tokens} (default 16{,}000),
the system prompt, the history, the tools, adaptive thinking and an effort
level (default \code{high}). It sends no sampling parameters, no forced tool
choice and no assistant prefill. Above the SDK's non-streaming output limit
the same request is streamed and the final message is read. A
\code{pause\_turn} stop is continued up to four times and the segments are
concatenated. SDK errors are classified from the most specific type to the
least. Server-side model fallback can change the model that answers, so it is
an explicit flag. It is off by default for benchmark runs, and the served
model of every segment is recorded, with refusal details and fallback
events.

\paragraph{Record and replay.}
\code{RecordingBackend} appends each turn to a JSONL file, keyed by the
SHA-256 of the canonical request: the backend configuration, the system
prompt, the full history and the tool definitions. A request already in the
recording is served from it, so a retried episode or a resumed run pays only
for turns it has not recorded; a recording is loaded, never truncated.
\code{ReplayBackend} serves recorded turns for identical requests and raises
an error on any miss. Tool results are deterministic functions of the data and
arguments, so a recorded language-model run can be re-scored offline, or under
a revised scorer.

\subsection{Audit log}
\label{sec:audit}

Each audit record is one line of canonical JSON with the fields \code{seq},
\code{ts}, \code{kind}, \code{episode\_id}, \code{payload}, \code{prev\_hash}
and \code{record\_hash}. The record hash is the SHA-256 of the record without
that field. The previous hash links to the preceding record, and the first
record links to 64 zeros. Verification checks the schema, the sequence
number, the link, the record hash, the canonical serialization and a trailing
newline. It therefore detects edited, deleted, inserted and reordered records,
and a partial final write. Removing records from the end of the file, or
rewriting the whole file with a recomputed chain, can only be detected against
an anchor stored elsewhere: the record count and head hash, which the runner
checks before writing any summary and records in the committed
\code{summary.json}. Scripted runs stamp records with a fixed time, so
rerunning them regenerates each log byte for byte and a regenerated log must
match the committed head. Opening an existing log verifies it first and
refuses to append to a broken chain. Keys that name
credentials are redacted, and so are API-key patterns and the current values
of credential environment variables.

\subsection{Grounding verifier}
\label{sec:grounding}

The verifier checks the numeric claims of an answer $A$ against the results
$R$ of its episode. Each result $r \in R$ has a set of numeric outputs $O_r$.

\paragraph{Claims.} Numbers are extracted as signed decimals, with thousands
separators, a leading \$, the units \%, \emph{percent}, \emph{pct},
percentage points, \emph{bp}, \emph{basis points} and multiples, and magnitude
suffixes (\emph{k}, \emph{M}, \emph{bn}, \emph{million}, \ldots). Unicode minus
signs and dashes are signs, and an unsigned number takes a verbal sign from an
adjacent cue such as ``down'', ``fell'' or ``a gain of''. The following are
excluded, and the reason is recorded: dates in ISO, slash and month-name
forms; clock times; years (four-digit unsigned integers from 1900 to 2100);
window lengths (an integer followed by a horizon noun, or preceded by
\emph{window}, \emph{lookback} or \emph{horizon}); function parameters such as
\code{vol(20)}; identifiers such as \code{SYN01}, \code{Q2} or \code{t+1};
ordinals and list markers; the constant in $\sqrt{252}$; numbers restated from
the question; and everything inside citation tags. A decimal glued to an
unknown suffix, or an implausibly long digit run, is an \emph{unparsed} claim
that counts and is never supported. Each remaining number is a claim $c$, with
value $v_c$, $d_c$ decimals shown, unit $u_c$, magnitude multiplier $m_c$, and
a flag $\sigma_c$ for an explicit or verbal sign.

\paragraph{Attribution.} Adjacent \code{[r:\dots]} tags form one citation
group. A claim is attributed to the first group after it in the same sentence.
If there is none, it is attributed to the nearest group before it in that
sentence; otherwise it is uncited. Sentences end at a newline or at a period,
exclamation or question mark followed by whitespace.

\paragraph{Rounding-aware match.} Let $S(\%) = S(\mathrm{pp}) = \{100\}$,
$S(\mathrm{bp}) = \{10^4\}$, $S(\$) = S(\times) = \{1\}$ and
$S(\text{plain}) = \{1, 100\}$, since tool outputs are fractions and models
often drop the percent sign. A claim $c$ matches an output $o$ if, for some
$s \in S(u_c)$,
\begin{equation}
  |x - v_c| \le \tfrac{1}{2}\,10^{-d_c}\,m_c + 10^{-9}\max(1, |v_c|),
  \qquad x = s\,o,
  \label{eq:match}
\end{equation}
and the signs agree: a signed claim must have the sign of $s\,o$, and an
unsigned claim matches only a non-negative $s\,o$, except for the maximum
drawdown, where $x = |s\,o|$ because ``a drawdown of 18.20\%'' is
conventional. The claim is then a correct rounding of the output under any tie
rule.

\paragraph{Binding.} Daily-bar results hold several fields and rows, so their
outputs are bound to the text. A row value such as the close of 2023-06-29
supports a claim only if that date appears in the claim's sentence or in the
question. A value with a bar field (open, high, low, close, volume, and the
range summaries) supports a claim only if the sentence, or else the question,
names no bar field or names that field. The first and last close of a range
are out of scope when the text speaks only of the other end. A claim that
matches only an out-of-scope output, or only an output of the opposite sign,
is unsupported and carries a note saying why.

\paragraph{Verdicts and rates.} An uncited claim is \emph{supported} if it
matches an output of any result in the episode. A cited claim is supported
only if it matches an output of a cited result that exists in the episode. A
claim whose group cites only identifiers absent from the episode is an
\emph{unknown citation} and counts as unsupported. The grounding rate is the
share of claims that are supported, and the citation rate is the share that
are cited. An episode is \emph{fully grounded} if it makes at least one claim
and all its claims are supported. The verifier checks faithfulness to tool
outputs, not whether the right quantity was computed. Section~\ref{sec:results}
shows a baseline that is fully grounded and still wrong.
